% Generated draft: requires mathematical and source review.
\documentclass[11pt]{article}
\usepackage[margin=25mm]{geometry}
\usepackage{fontspec}
\setmainfont{Latin Modern Roman}
\usepackage{amsmath,amssymb,mathtools,unicode-math}
\setmathfont{Latin Modern Math}
\usepackage{enumitem,needspace,xurl,hyperref}
\setlength{\parindent}{0pt}
\setlength{\parskip}{5pt}
\setlength{\emergencystretch}{3em}
\newcommand{\N}{\mathbb{N}}
\newcommand{\Z}{\mathbb{Z}}
\newcommand{\Q}{\mathbb{Q}}
\newcommand{\R}{\mathbb{R}}
\newcommand{\C}{\mathbb{C}}
\newcommand{\F}{\mathbb{F}}
\newcommand{\A}{\mathbb{A}}
\newcommand{\PP}{\mathbb{P}}
\newcommand{\E}{\mathbb{E}}
\newcommand{\eps}{\varepsilon}
\newcommand{\dd}{\,\mathrm{d}}
\begin{document}
\section*{UnsolvedMath ID: 30004725}
\textbf{Source problem code:} catalogue entry
\subsection{Definitions and mathematical statement}
Let \((X,g)\) be an asymptotically flat Riemannian manifold with compact boundary \(\partial X\), and let \(\operatorname{Scal}_g\) denote its scalar curvature. Assume
\[
\operatorname{Scal}_g(x)\geq 0\qquad\text{for every }x\in X.
\]
For a designated asymptotically flat end \(E\), write \(m_{\mathrm{ADM}}(E,g)\) for its ADM mass, provided the dimensional, decay, and integrability hypotheses needed to define that mass hold. The supplied record does not specify these hypotheses or designate an end when there is more than one.

A positive scalar-curvature lower bound on a nonempty region \(\Omega\subseteq X\) means that, for some constant \(\kappa>0\),
\[
\operatorname{Scal}_g(x)\geq\kappa\qquad\text{for every }x\in\Omega.
\]
When \(\partial X\neq\varnothing\), define the distance of this region from the boundary by
\[
d_g(\Omega,\partial X)
=\inf\{\operatorname{dist}_g(x,y):x\in\Omega,\ y\in\partial X\},
\]
where \(\operatorname{dist}_g\) is the Riemannian distance.

\textbf{Problem.} Find sufficient metric conditions involving such a positive lower bound \(\kappa\) on a region \(\Omega\) and its distance \(d_g(\Omega,\partial X)\) from the compact boundary that force
\[
m_{\mathrm{ADM}}(E,g)\geq 0.
\]
This is a request to determine sufficient conditions, rather than a specified implication with a supplied numerical threshold. The record gives no dimension, asymptotic decay convention, completeness assumption, boundary mean-curvature condition, admissible class of regions, or quantitative relation between curvature and distance. It also leaves unclear whether the conclusion concerns a single designated end or every end. Compactness of the boundary does not itself require the boundary to be nonempty; the intended treatment of an empty boundary is unspecified. No spin assumption appears in the question, and none is imposed here. These unresolved qualifications prevent the record from determining a single fully quantified conjecture.
\subsection{Short English statement}
For an asymptotically flat manifold with compact boundary and everywhere nonnegative scalar curvature, seek sufficient geometric conditions using a positive scalar-curvature bound on a region and that region's distance from the boundary to guarantee nonnegative ADM mass. The dimension, asymptotic hypotheses, region geometry, boundary conditions, and quantitative criterion remain unspecified. The supplied dated curation reports partial results in the spin setting and an unresolved general nonspin question; this assessment requires independent verification.
\subsection{Sources}
\begin{itemize}
\item Supplied local UnsolvedMath registry record, permanent ID \(30004725\), title: \emph{Localized Scalar Curvature Conditions Forcing Nonnegative ADM Mass}; exact source problem code: \texttt{catalogue entry}. Collection: Oberwolfach Reports Open Problems. The supplied proposed year is 2021; no individual proposer is named.
\item Supplied report citation: \emph{Analysis, Geometry and Topology of Positive Scalar Curvature Metrics} (2022), Oberwolfach Reports, pp.\ 1579--1649, DOI \texttt{10.4171/owr/2021/30}. \url{https://doi.org/10.4171/owr/2021/30}. The supplied background attributes the problem to workshop open-problem discussions documented by Oberwolfach Reports; no problem-specific page or speaker attribution is supplied.
\item Supplied literature source: ``Oberwolfach Reports 2022.'' \url{https://ems.press/content/serial-article-files/46909}.
\item Supplied literature source: ``Cecchini and Zeidler, curvature comparison.'' \url{https://arxiv.org/abs/2103.06833}.
\item Literature assessment dated exactly \texttt{2026-08-22}: the supplied curation labels the problem partially solved and reports explicit sufficient distance and scalar-curvature conditions in the spin setting, while reporting no known general nonspin characterization of the requested type. This is supplied dated curation derived from the OWR record, not an independently verified literature assessment; its status and cited claims require independent checking.
\item The registry labels the statement \texttt{corrected\_verified}, explaining that a publication-format audit preserved mathematical content while repairing notation or typography. This supplied audit label does not establish mathematical correctness or independent source verification for this draft.
\end{itemize}
\end{document}
